\documentclass[10pt,a4paper]{amsart}
\usepackage{amsmath,amssymb,amsthm,graphicx}
\usepackage{verbatim}
\newenvironment{reference}{\comment}{\endcomment}
%\newenvironment{reference}{}{}
\newenvironment{slogan}{\comment}{\endcomment}
\newenvironment{history}{\comment}{\endcomment}

% For commutative diagrams we use Xy-pic
\usepackage[all]{xy}

% We use 2cell for 2-commutative diagrams.
\xyoption{2cell}
% Theorem environments.
%
\theoremstyle{plain}
\newtheorem{theorem}[subsection]{Theorem}
\newtheorem{proposition}[subsection]{Proposition}
\newtheorem{lemma}[subsection]{Lemma}

\theoremstyle{definition}
\newtheorem{definition}[subsection]{Definition}
\newtheorem{example}[subsection]{Example}
\newtheorem{exercise}[subsection]{Exercise}
\newtheorem{situation}[subsection]{Situation}

\theoremstyle{remark}
\newtheorem{remark}[subsection]{Remark}
\newtheorem{remarks}[subsection]{Remarks}

\numberwithin{equation}{subsection}

% Macros
%
\def\lim{\mathop{\mathrm{lim}}\nolimits}
\def\colim{\mathop{\mathrm{colim}}\nolimits}
\def\Spec{\mathop{\mathrm{Spec}}}
\def\Hom{\mathop{\mathrm{Hom}}\nolimits}
\def\Ext{\mathop{\mathrm{Ext}}\nolimits}
\def\SheafHom{\mathop{\mathcal{H}\!\mathit{om}}\nolimits}
\def\SheafExt{\mathop{\mathcal{E}\!\mathit{xt}}\nolimits}
\def\Sch{\mathit{Sch}}
\def\Mor{\mathop{\mathrm{Mor}}\nolimits}
\def\Ob{\mathop{\mathrm{Ob}}\nolimits}
\def\Sh{\mathop{\mathit{Sh}}\nolimits}
\def\NL{\mathop{N\!L}\nolimits}
\def\CH{\mathop{\mathrm{CH}}\nolimits}
\def\proetale{{pro\text{-}\acute{e}tale}}
\def\etale{{\acute{e}tale}}
\def\QCoh{\mathit{QCoh}}
\def\Ker{\mathop{\mathrm{Ker}}}
\def\Im{\mathop{\mathrm{Im}}}
\def\Coker{\mathop{\mathrm{Coker}}}
\def\Coim{\mathop{\mathrm{Coim}}}

% Boxtimes
%
\DeclareMathSymbol{\boxtimes}{\mathbin}{AMSa}{"02}

%
% Macros for moduli stacks/spaces
%
\def\QCohstack{\mathcal{QC}\!\mathit{oh}}
\def\Cohstack{\mathcal{C}\!\mathit{oh}}
\def\Spacesstack{\mathcal{S}\!\mathit{paces}}
\def\Quotfunctor{\mathrm{Quot}}
\def\Hilbfunctor{\mathrm{Hilb}}
\def\Curvesstack{\mathcal{C}\!\mathit{urves}}
\def\Polarizedstack{\mathcal{P}\!\mathit{olarized}}
\def\Complexesstack{\mathcal{C}\!\mathit{omplexes}}
% \Pic is the operator that assigns to X its picard group, usage \Pic(X)
% \Picardstack_{X/B} denotes the Picard stack of X over B
% \Picardfunctor_{X/B} denotes the Picard functor of X over B
\def\Pic{\mathop{\mathrm{Pic}}\nolimits}
\def\Picardstack{\mathcal{P}\!\mathit{ic}}
\def\Picardfunctor{\mathrm{Pic}}
\def\Deformationcategory{\mathcal{D}\!\mathit{ef}}
\usepackage{cleveref}
\crefname{originalequation}{original equation}{original equations}
\Crefname{originalequation}{Original equation}{Original equations}
\makeatletter
\newlabel{equation-stalk}{{04EH}{source}{}{sourcepointer-1}{}}
\newlabel{equation-stalk@cref}{{[originalequation][0][]04EH}{[1][0][]source}{}{}{}}
\makeatother
\crefname{originallemma}{original lemma}{original lemmas}
\Crefname{originallemma}{Original lemma}{Original lemmas}
\makeatletter
\newlabel{categories-lemma-colimit-constant-connected-fibers}{{05US}{source}{}{sourcepointer-2}{}}
\newlabel{categories-lemma-colimit-constant-connected-fibers@cref}{{[originallemma][0][]05US}{[1][0][]source}{}{}{}}
\makeatother
\title{Stalks commute with the left Kan extension along a functor}
\author{The Stacks Project Authors}
\date{stacks-project:sites.tex:local-snapshot-20261012; GNU FDL 1.2 or later}
\begin{document}
\maketitle
\noindent\textit{Editorial note.} All delivered mathematical text is the author's own, copied from the pinned Stacks Project chapter file sites-cohomology.tex; no statement or proof step was rewritten, reordered or removed.\par\medskip
\noindent\textit{Editorial note.} The proof is the author's complete argument; every step the author writes out is inside the counted span and no closing sentence is a placeholder.\par\medskip
\noindent\textit{Editorial note.} Every labelled reference occurring inside the counted spans is carried in assembly.external\_references as a non-load-bearing navigation pointer (same-chapter ones at their definition site, cross-chapter ones at the invocation inside the counted span), so no load-bearing mathematics is replaced by a stub. No bibliography citation occurs inside the counted spans.\par\medskip
\noindent\textit{Editorial note.} The checked chapter list for same-name collisions is the set of stacks chapters staged in this repository: chow, cohomology, divisors, etale-cohomology, flat, more-algebra, more-morphisms, morphisms, spaces-more-morphisms, varieties together with this card's chapter sites-cohomology.\par\medskip
\noindent\textit{Editorial note.} The preamble is the chapter's own retained minimal preamble (sources/stacks-sites-cohomology/retained-preamble.tex), a verbatim excerpt of the native preamble: the Xy-pic setup and the native theorem declarations with the mathematical macros.\par\medskip
\noindent\textit{Allowed standard prerequisites (editorial).}\begin{itemize}
\item The description of the left Kan extension. Same chapter, the description of the left Kan extension as a colimit over morphisms into the images of objects. Supplies the first expression in the computation.
\item The description of the stalk functor. Same chapter, Equation equation-stalk. Supplies the description of the stalks appearing in the computation.
\item Colimits with connected fibres are unchanged. Categories, Lemma categories-lemma-colimit-constant-connected-fibers. Justifies the passage from the combined colimit to the colimit over the smaller index category.
\end{itemize}\medskip
\ifdefined\hypertarget\hypertarget{sourcepointer-1}{}\fi
\noindent\textit{Source reference.} Original equation 04EH: sources/stacks-sites/pinned-original-source.tex, source lines 8735--8735. Navigation pointer invoked at proof line 8735; the counted span contains the pointer only and the cited result is not reproduced inside it, so the pointer is registered as a non-load-bearing navigation stub in the established form for Stacks material. No bibliography citation is involved.\par\medskip
\ifdefined\hypertarget\hypertarget{sourcepointer-2}{}\fi
\noindent\textit{Source reference.} Original lemma 05US: sources/stacks-sites/pinned-original-source.tex, source lines 8754--8754. Navigation pointer invoked at proof line 8754; the counted span contains the pointer only and the cited result is not reproduced inside it, so the pointer is registered as a non-load-bearing navigation stub in the established form for Stacks material. No bibliography citation is involved.\par\medskip
\begin{lemma}
\label{lemma-point-functor}
Let $u : \mathcal{C} \to \mathcal{D}$ be a functor. Let
$v : \mathcal{D} \to \textit{Sets}$ be a functor and set
$w = v \circ u$. Denote $q$, resp., $p$ the stalk functor
(\ref{equation-stalk}) associated to $v$, resp.\ $w$.
Then $(u_p\mathcal{F})_q = \mathcal{F}_p$ functorially in the
presheaf $\mathcal{F}$ on $\mathcal{C}$.
\end{lemma}

\begin{proof}
This is a simple categorical fact. We have
\begin{align*}
(u_p\mathcal{F})_q
& =
\colim_{(V, y)} \colim_{U, \phi : V \to u(U)} \mathcal{F}(U) \\
& = \colim_{(V, y, U, \phi : V \to u(U))} \mathcal{F}(U) \\
& = \colim_{(U, x)} \mathcal{F}(U) \\
& = \mathcal{F}_p
\end{align*}
The first equality holds by the definition of $u_p$ and the
definition of the stalk functor. Observe that $y \in v(V)$.
In the second equality we simply combine colimits.
To see the third equality we apply
Categories, Lemma \ref{categories-lemma-colimit-constant-connected-fibers}
to the functor $F$ of diagram categories defined by the rule
$$
F((V, y, U, \phi : V \to u(U))) = (U, v(\phi)(y)).
$$
This makes sense because $w(U) = v(u(U))$. Let us check the hypotheses of
Categories, Lemma \ref{categories-lemma-colimit-constant-connected-fibers}.
Observe that $F$ has a right inverse, namely
$(U, x) \mapsto (u(U), x, U, \text{id} : u(U) \to u(U))$.
Again this makes sense because $x \in w(U) = v(u(U))$. On the other hand,
there is always a morphism
$$
(V, y, U, \phi : V \to u(U))
\longrightarrow
(u(U), v(\phi)(y), U, \text{id} : u(U) \to u(U))
$$
in the fibre category over $(U, x)$ which shows the fibre categories
are connected. The fourth and final equality is clear.
\end{proof}

\end{document}
